\section{Funzionalità aggiuntive}

Le funzionalità implementate nell'applicazione si dividono in due categorie principali: logico-funzionali e di interfaccia grafica (GUI).

Tra le funzionalità logico-funzionali rientrano:
\begin{itemize}
    \item Gestione di tre tipologie di attività differenti;
    \item Salvataggio e caricamento dei dati in formato JSON;
    \item Calcolo automatico dello stato dell'attività (in corso, completata, scaduta, ecc.) con comportamento personalizzato in base alla tipologia;
    \item Calcolo automatico della durata per la categoria degli impegni.
\end{itemize}

Per quanto riguarda le funzionalità legate all'interfaccia grafica, sono state implementate:
\begin{itemize}
    \item Suddivisione visiva automatica delle attività nelle rispettive colonne;
    \item Visualizzazione dettagliata di un'attività tramite doppio click;
    \item Sistema di sicurezza con avviso di emergenza in caso di chiusura con modifiche non salvate;
    \item Scorciatoie da tastiera per le azioni principali (salvataggio, caricamento, creazione ed eliminazione);
    \item Barra di ricerca reattiva in tempo reale e \textit{case-insensitive};
    \item Gestione dinamica del ridimensionamento della finestra;
    \item Alternanza dei colori nelle liste (\textit{zebra striping}), per facilitare la lettura e la distinzione visiva degli elementi.
\end{itemize}